\section{عدم الانعدام عند الواحد للشخصيات الحقيقية}

نغلق الآن الحالة الحقيقية من دون استعمال مبرهنة ديريشليه للأوليات في
المتتاليات الحسابية، ومن دون اللجوء إلى صيغة عدد الأصناف. الفكرة هي الجمع
بين إيجابية معاملات سلسلة ديريشليه مساعدة ومبدأ بسيط من نوع لاندو.

\begin{lemma}[الاستمرار التام وسلسلة ذات معاملات غير سالبة]
\label{lem:landau-positive-dirichlet-series}
\resultid{ANT-LEM-07-02}
\provedhere
لتكن
\[
F(s)=\sum_{n=1}^{\infty}\frac{a_n}{n^s},
\qquad a_n\geq0,
\]
متقاربة عندما \(\Re(s)>\sigma_0\). إذا امتدت \(F\) إلى دالة تامة، فإن
السلسلة تتقارب لكل \(s\in\mathbb R\).
\end{lemma}

\begin{proof}
ثبت عددًا حقيقيًا \(c>\sigma_0\)، ثم ثبت عددًا
\(\eta>0\) بحيث \(c-\eta>\sigma_0\). عندئذ، لكل \(k\geq0\)، لدينا
التقدير القياسي
\[
(\log n)^k=O_{k,\eta}(n^\eta),
\]
ومن ثم
\[
\sum_{n=1}^{\infty}
\frac{a_n(\log n)^k}{n^c}
\ll_{k,\eta}
\sum_{n=1}^{\infty}\frac{a_n}{n^{c-\eta}}<\infty.
\]
لذلك يجوز التفاضل حدًا حدًا في جوار \(c\)، فنحصل على
\[
F^{(k)}(c)
=
(-1)^k\sum_{n=1}^{\infty}
\frac{a_n(\log n)^k}{n^c}.
\]
ولأن \(F\) تامة، فإن متسلسلة تايلور حول \(c\) تتقارب عند كل عدد حقيقي
\(x<c\):
\[
F(x)
=
\sum_{k=0}^{\infty}\frac{F^{(k)}(c)}{k!}(x-c)^k.
\]
وبما أن \(a_n\geq0\) و\(c-x>0\)، تكون جميع الحدود بعد التعويض غير سالبة.
فتسمح مبرهنة تونيللي بتبديل الجمعين:
\begin{align*}
F(x)
&=
\sum_{k=0}^{\infty}\frac{(c-x)^k}{k!}
\sum_{n=1}^{\infty}\frac{a_n(\log n)^k}{n^c}\\
&=
\sum_{n=1}^{\infty}\frac{a_n}{n^c}
\sum_{k=0}^{\infty}
\frac{((c-x)\log n)^k}{k!}\\
&=
\sum_{n=1}^{\infty}\frac{a_n}{n^c}n^{c-x}
=
\sum_{n=1}^{\infty}\frac{a_n}{n^x}.
\end{align*}
وبما أن الطرف الأيسر منته، تتقارب السلسلة عند \(x\). أما \(x\geq c\)
فالتقارب معلوم أصلًا أو يتبع بالمقارنة.
\end{proof}

\begin{theorem}[عدم الانعدام للشخصيات الحقيقية]
\label{thm:real-character-nonvanishing-at-one}
\resultid{ANT-THM-07-11}
\provedhere
إذا كانت \(\chi\) شخصية ديريشليه حقيقية غير رئيسية، فإن
\[
L(1,\chi)\ne0.
\]
\end{theorem}

\begin{proof}
لنفترض عكس المطلوب، أي \(L(1,\chi)=0\). بما أن \(\chi\) غير رئيسية، فإن
\(L(s,\chi)\) دالة تامة. ودالة زيتا لها قطب بسيط وحيد عند \(s=1\).
فإذا كانت رتبة صفر \(L(s,\chi)\) عند \(1\) مساوية \(m\geq1\)، فإن
الجداء
\[
F(s)=\zeta(s)L(s,\chi)
\]
يمتد هولومورفيًا عبر \(s=1\): يكون غير صفري عند \(1\) إذا كان \(m=1\)،
ويملك صفرًا من الرتبة \(m-1\) إذا كان \(m>1\). ولا توجد لدالة زيتا
أقطاب أخرى، لذا تكون \(F\) دالة تامة.

عندما \(\Re(s)>1\)، يعطي ضرب سلسلتي ديريشليه
\[
F(s)
=
\sum_{n=1}^{\infty}\frac{a_n}{n^s},
\qquad
 a_n=\sum_{d\mid n}\chi(d).
\]
المعاملات \(a_n\) ضربية. وإذا كان \(p^k\) قوة أولية، فإن
\[
a_{p^k}=1+\chi(p)+\cdots+\chi(p)^k.
\]
ولأن \(\chi\) حقيقية، فـ\(\chi(p)\in\{-1,0,1\}\)، ولذلك
\[
a_{p^k}=
\begin{cases}
k+1,&\chi(p)=1,\\
1,&\chi(p)=0,\\
1,&\chi(p)=-1\text{ و }k\text{ زوجي},\\
0,&\chi(p)=-1\text{ و }k\text{ فردي}.
\end{cases}
\]
إذن \(a_n\geq0\) لكل \(n\).

أكثر من ذلك، لكل \(m\geq1\) لدينا \(a_{m^2}\geq1\)، لأن كل أس أولي في
\(m^2\) زوجي. ومن ثم
\[
\sum_{n=1}^{\infty}\frac{a_n}{n^{1/2}}
\geq
\sum_{m=1}^{\infty}\frac{a_{m^2}}{m}
\geq
\sum_{m=1}^{\infty}\frac1m
=\infty.
\]
لكن اللمّة السابقة تقول إن تمامية \(F\) مع عدم سلبية معاملاتها تجبر
سلسلتها على التقارب عند \(s=1/2\). هذا تناقض. لذلك
\(L(1,\chi)\ne0\).
\end{proof}

\begin{remark}
لا يعتمد هذا البرهان على وجود عدد أولي في كل فئة حسابية، ولا على صيغة عدد
الأصناف. وهو صالح للشخصيات الحقيقية غير الرئيسية سواء كانت بدائية أم
مستحثة، لأن الحجة تستعمل فقط تمامية \(L(s,\chi)\) وقيم
\(\chi(p)\in\{-1,0,1\}\).
\end{remark}

\section{إغلاق بوابة عدم الانعدام}

بضم المبرهنة \ref{thm:nonreal-character-nonvanishing-at-one} إلى المبرهنة
\ref{thm:real-character-nonvanishing-at-one} نحصل على عدم الانعدام لكل
شخصية غير رئيسية. وبذلك تغلق بوابة عدم الانعدام منطقيًا، ويبقى الانتقال
إلى مبرهنة ديريشليه للأوليات في المتتاليات الحسابية خطوة لاحقة مستقلة.
